\documentclass[11pt,letterpaper]{article}

\usepackage[utf8]{inputenc}
\usepackage[margin=1in]{geometry}
\usepackage{microtype}
\usepackage{parskip}
\usepackage{enumitem}
\usepackage{booktabs}
\usepackage{tabularx}
\usepackage{array}
\usepackage{amssymb}
\usepackage{framed}
\usepackage{listings}
\usepackage[hidelinks]{hyperref}

\hypersetup{pdftitle={Diaspore Build Guide}, pdfauthor={Rudr}}

\setcounter{secnumdepth}{1}
\setcounter{tocdepth}{2}
\setlength{\emergencystretch}{3em}

\newcommand{\code}[1]{\texttt{\detokenize{#1}}}
\newcommand{\lbl}[1]{\textbf{#1}\enspace}
\newcommand{\when}[1]{\par\textit{Planned: #1}\par}

\newlist{asks}{itemize}{1}
\setlist[asks]{label=\textbf{?}, leftmargin=1.6em, itemsep=2pt, topsep=2pt, parsep=0pt}
\newlist{checks}{itemize}{1}
\setlist[checks]{label=$\square$, leftmargin=1.6em, itemsep=2pt, topsep=2pt, parsep=0pt}
\setlist[itemize]{label=$\bullet$, leftmargin=1.6em, itemsep=2pt, topsep=2pt, parsep=0pt}
\setlist[enumerate]{leftmargin=1.8em, itemsep=2pt, topsep=2pt, parsep=0pt}

\newenvironment{milestone}[1]{\begin{framed}\noindent\textbf{#1}\par}{\end{framed}}
\newenvironment{tipbox}[1]{\begin{framed}\noindent\textbf{#1}\par}{\end{framed}}

\lstdefinestyle{json}{
  basicstyle=\ttfamily\small,
  breaklines=true,
  columns=fullflexible,
  keepspaces=true,
  frame=single,
  framesep=4pt,
  xleftmargin=4pt,
  xrightmargin=4pt,
  showstringspaces=false
}

\title{\textbf{Diaspore Build Guide}\\[0.3em]\large Checkpoints from September 23 to December 14, 2026}
\author{Rudr}
\date{Version 0.1, revised October 7, 2026}

\begin{document}

\maketitle

\tableofcontents

\newpage

\section{How to use this guide}

This guide splits Diaspore into small checkpoints. Each one teaches one idea and ends with working, tested code.

Every checkpoint has the same shape:

\begin{itemize}
  \item \lbl{Goal} What works at the end.
  \item \lbl{Why} What breaks without it.
  \item \lbl{Learn} The one new idea, in plain words.
  \item \lbl{Decide first} Questions to answer on paper before coding.
  \item \lbl{Tests first} Behaviors to write tests for before writing the code.
  \item \lbl{Done when} How you know it is time to move on.
  \item \lbl{Go you will meet} Packages and features you will likely use.
\end{itemize}

The rules:

\begin{enumerate}
  \item One checkpoint at a time. Do not start the next one until every ``Done when'' box is checked.
  \item Write your ``Decide first'' answers in the decisions log (Section~\ref{sec:log}) before coding.
  \item Tests before code, every time.
  \item No function signatures are given. Designing your own types is part of the learning.
  \item If you read someone else's code for an idea, close it before writing yours. Then write yours from memory.
\end{enumerate}

\begin{tipbox}{When you are stuck}
\begin{enumerate}
  \item Name the data shape in one plain sentence.
  \item Write the concrete, repetitive version first, even if it looks clumsy.
  \item Replace the parts that change with a variable.
  \item Draw it: a memory diagram of what exists before and after one step.
\end{enumerate}
\end{tipbox}

\section{What Diaspore is}

Diaspore is a deterministic simulator for distributed systems. It is written in Go and ships as one binary.

Nodes are separate programs, written in any language. They talk to Diaspore through JSON lines on stdin and stdout. SDKs are optional helpers, not requirements.

Diaspore owns time, the network, crashes, and randomness, all driven by one seed. It sends one event to one node at a time and waits for the reply. That is what makes every run repeat exactly.

When a rule breaks, Diaspore saves a \code{.pappus} file. Anyone can replay it and see the same bug.

The commands are \code{diaspore run}, \code{diaspore dandelion}, \code{diaspore replay}, and \code{diaspore check-determinism}.

\lbl{Honest pitch} ``A small, open, language-agnostic take on FoundationDB-style simulation, inspired by FoundationDB, TigerBeetle, and Maelstrom.'' Never ``the first.''

\subsection{Real code, simulated world}

Think of a flight simulator. The pilot is real, but the sky is simulated. Lessons from a simulated crash still hold, because the pilot makes the same decisions in a real plane.

Diaspore works the same way. The bug lives in a system's logic, and that logic can run inside Diaspore even though the world around it is fake.

Sans-I/O literally means ``without input and output.'' The core logic never touches the network or the clock directly. It is only handed events and answers them.

Code built this way runs in Diaspore for testing and in production for real use. Only a thin layer differs: in production it talks to real sockets and clocks, and under Diaspore it talks JSON lines.

etcd's Raft library is built this way. FoundationDB and TigerBeetle test their real production code like this too.

The honest cost: a system must be built, or adapted, in this style first.

\subsection{Where it is useful}

\begin{itemize}
  \item \textbf{Building replicated systems.} Leader election, replication, and consensus code hit bugs that need a crash plus bad timing. A sweep over thousands of seeds finds them.
  \item \textbf{Bug reports.} Instead of ``it fails sometimes,'' a developer attaches a \code{.pappus} file. Anyone can replay the exact failure on their own machine.
  \item \textbf{Regression tests in CI.} Once a bug is fixed, its \code{.pappus} file becomes a test that proves it stays fixed. Bonus B3 runs a sweep on every pull request.
  \item \textbf{Learning distributed systems.} Students building Raft or a key-value store can see exactly why their code broke. Maelstrom already serves this group, so exact replay is the edge.
\end{itemize}

\subsection{What it can and cannot control}

Diaspore controls what travels between nodes, and when things happen. It cannot control what happens inside the machine: threads, disk, memory, the kernel, or containers.

\begin{center}
\small
\begin{tabularx}{\linewidth}{@{}X l l@{}}
\toprule
\textbf{Issue} & \textbf{v0.1, florets} & \textbf{v0.2, interception} \\
\midrule
A message arrives late or out of order & Yes & Yes \\
A message is lost, or a partition cuts nodes off & Yes & Yes \\
A message arrives twice & Yes & Yes \\
A node crashes and restarts & Yes & Partly, since disk is real \\
A timeout fires before a slow reply arrives & Yes & Yes \\
Two nodes' clocks disagree & Not yet & Possible \\
Two threads inside one node race each other & No & No \\
A disk write is lost before it is saved & No & No \\
Docker networking or DNS is misconfigured & No & No \\
A node runs out of memory or CPU & No & No \\
Real TCP or kernel behavior misbehaves & No & No \\
\bottomrule
\end{tabularx}
\end{center}

Here v0.1 and v0.2 are Diaspore releases, not versions of this guide.

\section{Scope}

\subsection{Must-have: the course version}

\begin{itemize}
  \item Core: event loop, virtual time, one seed, and a trace hash.
  \item Protocol v1, written as a spec with JSON Schema files.
  \item A lockstep process driver with a watchdog.
  \item Timers, storage, and crash with restart.
  \item Faults: delay, drop, duplicate, crash, and partition.
  \item Workload and history, plus two checkers: acked writes and linearizability.
  \item A raw Go key-value example node with one planted bug.
  \item One tiny raw node in another language.
  \item CLI: \code{run}, \code{dandelion} across all cores, \code{replay}, and \code{check-determinism}.
  \item The \code{.pappus} v1 format.
  \item Scaling measurements for the report.
\end{itemize}

\subsection{Bonus, in this order}

\begin{enumerate}
  \item \code{dandelion} across machines.
  \item An HTML timeline viewer.
  \item A GitHub Action.
  \item A Go SDK helper.
  \item etcd's Raft library, wrapped as a floret.
\end{enumerate}

\subsection{Not now}

Python SDK, a real production runtime, disk faults, clock skew, shrinking failing runs, and binary encodings like protobuf.

\subsection{After the course: new failure classes}

These fit the building-block design, so you or the community can add them after v0.1.

\begin{itemize}
  \item \textbf{Gray failure} means a floret is degraded but still looks healthy. Diaspore would create the conditions, like a slow floret, a flaky link, or a one-way partition, and the checkers would catch the damage.
  \item \textbf{Response failure} means a floret returns wrong values or makes an invalid state change. Diaspore would change a reply in transit, or hand back corrupted stored state on restart.
  \item \textbf{Byzantine failure} means a floret lies, or tells different peers different things. Diaspore would change a floret's messages differently for each recipient. Realistic lies need knowledge of the message format, so this suits community plugins.
\end{itemize}

In every case the seed decides the change, so runs stay deterministic. Tolerating Byzantine faults, with signing and special quorums, is the tested system's job, not Diaspore's.

\subsection{After the course: interception mode}

The long-term goal is for Diaspore to wrap unmodified programs, not only florets. Interception literally means catching a program's requests to the OS, like ``what time is it'' or ``send this packet,'' and answering them yourself.

\begin{itemize}
  \item Diaspore would answer every clock read, random number, and network call, so it dictates the program's whole environment.
  \item It starts Linux only, with single-threaded programs.
  \item Real programs ask questions in the middle of their work, so the event model needs a new kind of call: the program asks, and waits for Diaspore's answer.
  \item Go programs skip the C library, so \code{LD_PRELOAD} tricks do not work on them.
  \item Linux often answers clock reads without a real system call, so a simple interceptor never sees them.
\end{itemize}

This is why the loop talks to florets only through a small driver interface, decided in CP9. Interception becomes a second driver beside the process driver, and the core stays the same.

\section{Calendar}

Assumptions:

\begin{itemize}
  \item About 8 to 10 hours a week go to Diaspore. Each checkpoint is sized for 2 to 4 hours.
  \item Course work is due Mondays at 9am Vancouver time, so keep Sundays for assignments.
  \item Thanksgiving is Thursday, November 26. That week is planned light.
  \item The semester ends December 14. Confirm the exact final project due date and presentation date on the syllabus.
\end{itemize}

\begin{center}
\small
\begin{tabularx}{\linewidth}{@{}l l l X l@{}}
\toprule
\textbf{Week} & \textbf{Dates} & \textbf{Phase} & \textbf{Checkpoints} & \textbf{Milestone} \\
\midrule
1  & Sep 23--27     & Ground     & CP0                  & \\
2  & Sep 28--Oct 4  & Core       & CP1--CP3             & \\
3  & Oct 5--11      & Core       & CP4--CP6             & A: Oct 11 \\
4  & Oct 12--18     & Protocol   & CP7--CP9             & \\
5  & Oct 19--25     & Protocol   & CP10--CP12           & B: Oct 25, Gate 1 \\
6  & Oct 26--Nov 1  & First bug  & CP13--CP14           & \\
7  & Nov 2--8       & First bug  & CP15--CP16           & \\
8  & Nov 9--15      & First bug  & CP17--CP19           & C: Nov 15 \\
9  & Nov 16--22     & Finish     & CP20--CP21           & \\
10 & Nov 23--29     & Finish     & CP22--CP23 (light week) & D: Nov 29, Gate 2 \\
11 & Nov 30--Dec 6  & Bonus      & B1, then B2          & \\
12 & Dec 7--13      & Wrap-up    & Report, demo, buffer & Ends Dec 14 \\
\bottomrule
\end{tabularx}
\end{center}

\lbl{Gate 1, Sunday, October 25} If Milestone B is not working, drop every bonus item. Keep going on must-haves only.

\lbl{Gate 2, Sunday, November 29} If Milestone D is not done, skip the bonus week. Use it to finish the must-haves.

\section{Phase 0: Ground}

\subsection{CP0: A repo that tests itself}
\when{September 23 to 27}

\lbl{Goal} An empty Go project where every push runs the tests.

\lbl{Why} Without it, you find broken code days later.

\lbl{Learn} CI, short for continuous integration, literally means a server that runs your checks on every push.

\lbl{Decide first}
\begin{asks}
  \item What are the module path and the repo name?
  \item Apache-2.0 or MIT? Apache-2.0 adds a patent grant. Both are common for developer tools.
  \item Read the repository layout in \code{docs/design.md}. Can you say what each folder is for, in one sentence each?
\end{asks}

\lbl{Tests first}
\begin{checks}
  \item One trivial test that passes.
  \item One test you break on purpose, to watch CI fail.
\end{checks}

\lbl{Done when}
\begin{checks}
  \item CI goes red on a broken test and green after the fix.
  \item CI runs on Linux, macOS, and Windows.
  \item The README has the one-line pitch.
\end{checks}

\lbl{Go you will meet} \code{go mod}, \code{go test}, \code{go vet}, and a GitHub Actions workflow file.

\section{Phase 1: The deterministic core}

Everything in this phase runs inside one Go process. There is no JSON and no child process yet.

One rule holds for the whole phase: no goroutines and no real clock in the core.

\subsection{CP1: Virtual clock}
\when{Week of September 28}

\lbl{Goal} A clock that moves only when the simulator moves it.

\lbl{Why} Real time is different on every run, so replays could never match.

\lbl{Learn} Virtual time literally means time is just a number the simulator owns. Ten simulated minutes can pass in one real second.

\lbl{Decide first}
\begin{asks}
  \item Say the clock's data shape in one sentence.
  \item Milliseconds or microseconds? A plain integer or a duration type? Write down why.
  \item Who may move the clock? What happens if someone tries to move it backward?
\end{asks}

\lbl{Tests first}
\begin{checks}
  \item The clock starts at zero.
  \item Moving forward works.
  \item Moving backward is refused.
\end{checks}

\lbl{Done when}
\begin{checks}
  \item All tests pass.
  \item Nothing in the package reads the real clock.
\end{checks}

\lbl{Go you will meet} \code{time.Duration}, \code{int64}, and the \code{errors} package.

\subsection{CP2: Event queue}
\when{Week of September 28}

\lbl{Goal} A queue that always hands back the earliest event.

\lbl{Why} Events are created out of order, but they must happen in time order.

\lbl{Learn} A priority queue literally means a line where the most urgent item always comes out first. Here, most urgent means earliest.

\lbl{Decide first}
\begin{asks}
  \item What does every event need to carry? Say it in one sentence.
  \item Two events at the same time: which goes first? Why must that answer never depend on luck?
  \item Values or pointers in the queue? Draw both in memory before choosing.
\end{asks}

\lbl{Tests first}
\begin{checks}
  \item Events come out in time order.
  \item Same-time events come out in the order they were added.
  \item Taking from an empty queue behaves the way you decided.
  \item 10,000 inserts made from a fixed seed come out sorted.
\end{checks}

\lbl{Done when}
\begin{checks}
  \item All tests pass.
  \item You can draw the heap after five inserts, by hand.
\end{checks}

\lbl{Go you will meet} \code{container/heap}.

\subsection{CP3: One seed, many streams}
\when{Week of September 28}

\lbl{Goal} Every random choice in a run comes from one seed.

\lbl{Why} One stray random call anywhere breaks replay.

\lbl{Learn} A PRNG, or pseudo-random number generator, literally means a formula that makes numbers look random but repeats exactly from the same seed.

\lbl{Decide first}
\begin{asks}
  \item One generator for everything, or one child stream per purpose, like network, faults, workload, and each node?
  \item If you add a new fault type later, should every old seed's network delays change? Let your answer pick the design.
  \item How do you turn the master seed plus a name like ``network'' into a child seed, the same way on every machine?
\end{asks}

\lbl{Tests first}
\begin{checks}
  \item The same seed gives the same first 1,000 numbers.
  \item Different seeds give different numbers.
  \item Adding a new child stream does not change the existing ones.
\end{checks}

\lbl{Done when}
\begin{checks}
  \item All tests pass.
  \item No global random functions are used anywhere. Check with grep.
\end{checks}

\lbl{Go you will meet} \code{math/rand/v2}, whose PCG generator takes explicit seeds, and \code{hash/fnv}.

\subsection{CP4: The loop}
\when{Week of October 5}

\lbl{Goal} A loop that takes the next event, runs its handler, and repeats until the run ends.

\lbl{Why} This is the heart of Diaspore. Everything else plugs into it.

\lbl{Learn} A discrete-event simulation literally means jumping from one event straight to the next, skipping the empty time between them.

\lbl{Decide first}
\begin{asks}
  \item What does a handler receive, and what may it do? Can it call another node directly, or only schedule events?
  \item When does a run stop? List every reason.
  \item Where do test nodes live for now?
\end{asks}

\lbl{Tests first}
\begin{checks}
  \item Two fake nodes play ping-pong ten times.
  \item The final virtual time equals the sum of the delays you set.
  \item Scheduling an event in the past is refused.
  \item The run stops at its time limit.
\end{checks}

\lbl{Done when}
\begin{checks}
  \item All tests pass.
  \item The loop has no goroutines.
\end{checks}

\subsection{CP5: Trace and hash}
\when{Week of October 5}

\lbl{Goal} Every run writes a trace, and the trace gets a fingerprint.

\lbl{Why} You need proof that two runs were exactly the same.

\lbl{Learn} A hash literally means a short fingerprint of some data. Change one byte and the fingerprint changes.

\lbl{Decide first}
\begin{asks}
  \item What goes in one trace line? Enough to debug, and nothing more.
  \item Will the same event always turn into the same bytes? Think about map key order and number formatting.
  \item Hash while writing, or after the run?
\end{asks}

\lbl{Tests first}
\begin{checks}
  \item The same seed twice gives the same hash.
  \item Changing one delay changes the hash.
  \item Every trace line parses back.
\end{checks}

\lbl{Done when}
\begin{checks}
  \item The ping-pong test also checks the hash.
\end{checks}

\lbl{Go you will meet} \code{encoding/json}, \code{crypto/sha256}, \code{io.MultiWriter}, and \code{bufio}.

\subsection{CP6: Simulated network}
\when{Week of October 5}

\lbl{Goal} Messages between nodes travel through Diaspore with a seeded delay.

\lbl{Why} The network causes most distributed bugs, so Diaspore must control it.

\lbl{Decide first}
\begin{asks}
  \item What does the network hold? Say it in one sentence.
  \item Where does the delay range come from for now?
  \item Random delays already let messages overtake each other. Do you still need a separate reorder fault? Write your answer in the log.
\end{asks}

\lbl{Tests first}
\begin{checks}
  \item Every message arrives inside its delay range.
  \item The same seed gives the same arrival times.
  \item With a wide delay range, some messages overtake others.
\end{checks}

\lbl{Done when}
\begin{checks}
  \item Ping-pong runs over the simulated network.
\end{checks}

\begin{milestone}{Milestone A --- Sunday, October 11: Deterministic core}
\begin{checks}
  \item Ping-pong runs over the simulated network.
  \item The same seed gives the same hash, every time.
  \item The core has no goroutines and no real clock.
  \item Optional: draft blog post 1, on why distributed bugs are hard to replay.
\end{checks}
\end{milestone}

\section{Phase 2: Real processes and the protocol}

Now nodes become real programs. The protocol is the product, so it gets written down first.

\subsection{CP7: Protocol spec v1}
\when{Week of October 12}

\lbl{Goal} A written spec that someone could implement without asking you anything.

\lbl{Why} SDKs are optional, so the spec is what people build against.

\lbl{Learn} JSON Schema literally means a JSON file that describes what valid JSON looks like.

\lbl{Decide first}
\begin{asks}
  \item Start from the draft in Appendix~\ref{app:protocol}. What would you change?
  \item How does a node learn the protocol version? What happens on a mismatch?
  \item Which reply fields are required? What should a node do with an event type it does not know?
\end{asks}

\lbl{Tests first}
\begin{checks}
  \item The schema accepts every example in Appendix~\ref{app:protocol}.
  \item The schema rejects five broken examples that you write.
\end{checks}

\lbl{Done when}
\begin{checks}
  \item \code{PROTOCOL.md} and the schema files are in the repo.
\end{checks}

\lbl{Go you will meet} A JSON Schema validator library for your tests, if you want one.

\subsection{CP8: Talking to one process}
\when{Week of October 12}

\lbl{Goal} Start a node program, send it \code{init}, and read its reply.

\lbl{Why} Nodes are separate programs, so Diaspore drives them through pipes.

\lbl{Learn} A pipe literally means a one-way stream of bytes between two programs. A program reads from its stdin and writes to its stdout.

\lbl{Decide first}
\begin{asks}
  \item Who owns the process and its pipes? What gets closed, and when?
  \item How do you read exactly one JSON line? What if a line is very long?
  \item Where do node logs go, so stdout stays clean for the protocol?
\end{asks}

\lbl{Tests first}
\begin{checks}
  \item A tiny raw test node, written in Go with no SDK, answers \code{init}.
  \item A node that prints garbage gives a clear error.
  \item A node that exits early gives a clear error.
\end{checks}

\lbl{Done when}
\begin{checks}
  \item All tests pass on all three operating systems in CI.
\end{checks}

\lbl{Go you will meet} \code{os/exec}, \code{bufio.Scanner} (note its default line limit of 64 KB), \code{bufio.Writer} with \code{Flush}, and \code{encoding/json}.

\subsection{CP9: Lockstep driver}
\when{Week of October 12}

\lbl{Goal} Send one event, wait for its reply, apply it, and repeat.

\lbl{Why} Without lockstep, processes race each other and runs stop matching.

\lbl{Learn} Lockstep literally means moving one step at a time, together, never ahead. A seam literally means a boundary where one piece can be swapped without touching the rest. In Go, a seam is usually an interface.

\lbl{Decide first}
\begin{asks}
  \item What happens to the \code{send} list in a reply? Say it in one sentence.
  \item A reply has sends, timers, and storage. Does the order you apply them in matter?
  \item How long does the watchdog wait in real time? What does a timeout mean for the run?
  \item What is the smallest interface between the loop and a floret driver? Keep the process driver behind it, so a second driver, like interception, can be added later without changing the loop.
\end{asks}

\lbl{Tests first}
\begin{checks}
  \item Three test nodes pass a message around a ring.
  \item Twenty runs of one seed give one hash.
  \item A node that never replies trips the watchdog with a clear message.
\end{checks}

\lbl{Done when}
\begin{checks}
  \item All tests pass.
  \item The watchdog can stop a run, but it never changes the run's results.
\end{checks}

\lbl{Go you will meet} Goroutines and \code{select}, but only in the process driver, never in the loop. Also \code{context.WithTimeout}.

\subsection{CP10: Timers}
\when{Week of October 19}

\lbl{Goal} Nodes can set and cancel timers.

\lbl{Why} Heartbeats, retries, and election timeouts all need timers.

\lbl{Decide first}
\begin{asks}
  \item Setting a timer whose id already exists: replace it, or refuse?
  \item Canceling a timer that already fired: ignore it, or refuse?
  \item A timer and a message at the same instant: which goes first? Your CP2 tie-break should answer this.
\end{asks}

\lbl{Tests first}
\begin{checks}
  \item A timer fires exactly at the current time plus its delay.
  \item Canceling stops it from firing.
  \item Setting it again moves it.
  \item A node with a 100 ms heartbeat, run for one second, sends the number of heartbeats you predicted on paper.
\end{checks}

\lbl{Done when}
\begin{checks}
  \item All tests pass.
\end{checks}

\subsection{CP11: Storage, crash, and restart}
\when{Week of October 19}

\lbl{Goal} Nodes can save state that survives a crash.

\lbl{Why} Many real bugs appear only after a crash and a restart.

\lbl{Learn} Durable literally means it survives a crash. Memory does not.

\lbl{Decide first}
\begin{asks}
  \item What does Diaspore keep for each node's storage? One sentence.
  \item Does \code{store} merge keys, or replace everything?
  \item On a crash, what happens to messages already on their way to that node? To its timers?
  \item On restart, start a fresh process or reuse the old one? Which is more honest?
\end{asks}

\lbl{Tests first}
\begin{checks}
  \item A counter node crashes, restarts, and continues from its stored value.
  \item Messages and timers behave the way you decided.
  \item Crash and restart show up in the trace.
\end{checks}

\lbl{Done when}
\begin{checks}
  \item All tests pass.
\end{checks}

\subsection{CP12: check-determinism}
\when{Week of October 19}

\lbl{Goal} A command that runs one seed twice and compares the results.

\lbl{Why} In any language, a node can break determinism by accident. This catches it.

\lbl{Decide first}
\begin{asks}
  \item Report only ``different,'' or point to the first line that differs? Which helps more?
  \item Which exit code means ``not deterministic''?
  \item The \code{flag} package or a CLI library? Write down why.
\end{asks}

\lbl{Tests first}
\begin{checks}
  \item A well-behaved node passes.
  \item A node that reads the real clock fails, and the report points at the first differing line.
\end{checks}

\lbl{Done when}
\begin{checks}
  \item \code{diaspore check-determinism --seed 42} works on the ring test.
\end{checks}

\begin{milestone}{Milestone B --- Sunday, October 25: Real processes in lockstep (Gate 1)}
\begin{checks}
  \item Real node processes run in lockstep, with timers.
  \item Crash and restart work, with storage.
  \item \code{check-determinism} catches a bad node.
  \item Publish blog post 1.
\end{checks}
If this is not working, drop every bonus item now.
\end{milestone}

\section{Phase 3: The first real bug}

By the end of this phase, a sweep finds the planted bug and a \code{.pappus} file replays it every time.

\subsection{CP13: Config file}
\when{Week of October 26}

\lbl{Goal} \code{diaspore.yaml} describes a whole run.

\lbl{Why} Runs must be shareable and repeatable without long command lines.

\lbl{Decide first}
\begin{asks}
  \item Which fields are required, and what are the defaults?
  \item How do you report a bad field so the user can fix it fast?
  \item Which YAML library? Check that it is still maintained before picking.
\end{asks}

\lbl{Tests first}
\begin{checks}
  \item The example in Appendix~\ref{app:files} loads.
  \item Unknown fields are rejected.
  \item A bad duration like ``5 seconds'' is rejected with a clear message.
\end{checks}

\lbl{Done when}
\begin{checks}
  \item All tests pass.
  \item \code{diaspore run} reads everything it needs from the file.
\end{checks}

\subsection{CP14: Faults: delay, drop, and duplicate}
\when{Week of October 26}

\lbl{Goal} Rules from the config change what the network does.

\lbl{Why} Bugs hide behind unlucky delivery, so Diaspore must create that bad luck on purpose.

\lbl{Learn} A fault rule literally means three things: which messages, what goes wrong, and how often.

\lbl{Decide first}
\begin{asks}
  \item Say a rule's shape in one sentence.
  \item If two rules match one message, do both apply? In what fixed order?
  \item Which random stream do faults draw from?
\end{asks}

\lbl{Tests first}
\begin{checks}
  \item Drop with chance 1 drops every matching message. Chance 0 drops none.
  \item Duplicate delivers a message twice.
  \item Rules touch only the messages they match.
  \item The same seed gives the same drops.
\end{checks}

\lbl{Done when}
\begin{checks}
  \item Every fault shows up in the trace, along with the rule that caused it.
\end{checks}

\subsection{CP15: Crash schedule and partitions}
\when{Week of November 2}

\lbl{Goal} Crash nodes on a seeded schedule, and split the network into groups.

\lbl{Why} Crashes and partitions are where replication bugs live.

\lbl{Learn} A network partition literally means some nodes cannot reach others for a while.

\lbl{Decide first}
\begin{asks}
  \item Say a partition's shape in one sentence.
  \item Messages that cross a partition: dropped, or held until it heals?
  \item How are crash and partition times picked from the seed?
\end{asks}

\lbl{Tests first}
\begin{checks}
  \item During a partition, no message crosses between groups.
  \item After it heals, messages flow again.
  \item The same seed gives the same schedule.
\end{checks}

\lbl{Done when}
\begin{checks}
  \item All tests pass.
\end{checks}

\subsection{CP16: Workload and history}
\when{Week of November 2}

\lbl{Goal} Virtual clients send key-value requests and record what happened.

\lbl{Why} Checkers can only judge what was recorded.

\lbl{Learn} A history literally means the list of what every client asked, what it got back, and when.

\lbl{Decide first}
\begin{asks}
  \item What is one history entry? One sentence.
  \item When does a client give up waiting, and what do you record then? A timed-out write may or may not have happened.
  \item How do clients pick keys and values from the seed?
\end{asks}

\lbl{Tests first}
\begin{checks}
  \item There is one history entry per request.
  \item Timeouts are recorded as unknown, not as failed.
  \item The same seed gives the same history.
\end{checks}

\lbl{Done when}
\begin{checks}
  \item \code{history.jsonl} is written for every run.
\end{checks}

\subsection{CP17: Example key-value node with a planted bug}
\when{Week of November 9}

\lbl{Goal} A small replicated key-value node in Go, speaking raw JSON with no SDK.

\lbl{Why} Diaspore needs a real target, and a known bug proves the tool works.

\lbl{Decide first}
\begin{asks}
  \item Describe the replication design in three sentences.
  \item Pick the planted bug. Choose one that only shows up under a crash plus unlucky timing.
  \item How does each node know which one is the primary?
\end{asks}

\lbl{Tests first}
\begin{checks}
  \item With no faults, every write can be read back.
  \item The node passes \code{check-determinism}.
\end{checks}

\lbl{Done when}
\begin{checks}
  \item Fault-free runs are clean.
\end{checks}

\subsection{CP18: Checker: acked writes are never lost}
\when{Week of November 9}

\lbl{Goal} A checker that catches the planted bug.

\lbl{Why} Without a rule to check, a broken run looks like any other run.

\lbl{Learn} An invariant literally means a rule that must always stay true.

\lbl{Decide first}
\begin{asks}
  \item Write the rule in one exact sentence.
  \item How do unknown outcomes count?
  \item Check once at the end of the run, or after every read?
\end{asks}

\lbl{Tests first} Use hand-written histories:
\begin{checks}
  \item A clean history passes.
  \item A history with a lost acked write fails.
  \item A history with an unknown write does not fail by mistake.
\end{checks}

\lbl{Done when}
\begin{checks}
  \item Clean runs pass, and the planted bug fails.
\end{checks}

\subsection{CP19: .pappus, replay, and the first sweep}
\when{Week of November 9}

\lbl{Goal} Sweep seeds until the bug appears, save it, and replay it.

\lbl{Why} This is the whole promise of Diaspore, working end to end.

\lbl{Decide first}
\begin{asks}
  \item What must a \code{.pappus} file hold to stand alone? Start from Appendix~\ref{app:files}.
  \item If the node program changed since the file was made: warn, or refuse?
  \item Should a sweep stop at the first failure, or collect all of them?
\end{asks}

\lbl{Tests first}
\begin{checks}
  \item Writing and then reading a \code{.pappus} file gives back the same content.
  \item Replay gives the same trace hash as the original run.
  \item Replay warns when the node program changed.
\end{checks}

\lbl{Done when}
\begin{checks}
  \item \code{diaspore dandelion --seeds 1-500} finds the bug.
  \item \code{diaspore replay} shows it every time.
\end{checks}

\begin{milestone}{Milestone C --- Sunday, November 15: The first bug, end to end}
\begin{checks}
  \item A sweep finds the planted bug.
  \item The \code{.pappus} file replays it with the same hash.
  \item Record a short demo: sweep, failure, replay.
  \item Optional: blog post 2, on finding the first bug.
\end{checks}
\end{milestone}

\section{Phase 4: Finish the must-haves}

\subsection{CP20: Parallel dandelion and scaling numbers}
\when{Week of November 16}

\lbl{Goal} Run many seeds at once across CPU cores, and measure how it scales.

\lbl{Why} This is the scalability half of the course story.

\lbl{Learn} Scaling up literally means using more of one machine. Scaling out means using more machines, which is bonus B1.

\lbl{Decide first}
\begin{asks}
  \item Parallelism lives between seeds, never inside one run. Why?
  \item How many workers? How are results collected? How do all workers stop at the first failure?
  \item What will you measure? At least: seeds per second against worker count, events per second per run, and time spent waiting on node processes.
\end{asks}

\lbl{Tests first}
\begin{checks}
  \item A parallel sweep finds the same failing seeds as a serial sweep.
  \item Results do not depend on the number of workers.
\end{checks}

\lbl{Done when}
\begin{checks}
  \item You have a chart of seeds per second for 1, 2, 4, and 8 workers.
\end{checks}

\lbl{Go you will meet} Goroutines, channels, \code{sync.WaitGroup}, context cancellation, and \code{runtime.NumCPU}.

\subsection{CP21: Linearizability checker with Porcupine}
\when{Week of November 16}

\lbl{Goal} A stronger checker for the key-value store.

\lbl{Why} ``No lost writes'' misses bugs like stale reads.

\lbl{Learn} Linearizability literally means every operation looks like it happened at a single instant between its start and its end, as if there were only one copy of the data.

\lbl{Decide first}
\begin{asks}
  \item Read Porcupine's README first. How will you describe the key-value store as a model?
  \item How do unknown outcomes map onto that model?
  \item Check all keys together, or one key at a time?
\end{asks}

\lbl{Tests first}
\begin{checks}
  \item A hand-made valid history passes.
  \item A history with a stale read fails.
\end{checks}

\lbl{Done when}
\begin{checks}
  \item The checker can be turned on by name in \code{diaspore.yaml}.
\end{checks}

\lbl{Go you will meet} \code{github.com/anishathalye/porcupine}.

\subsection{CP22: One tiny raw node in another language}
\when{Week of November 23 (light week)}

\lbl{Goal} Prove ``any language'' with a small node written outside Go.

\lbl{Why} Without it, the honest wording is only ``designed for any language.''

\lbl{Decide first}
\begin{asks}
  \item JavaScript or Python?
  \item What are that language's determinism traps? Python randomizes string hashing per process. JavaScript's \code{Math.random} cannot be seeded.
  \item How will Diaspore set up each node's environment to avoid those traps?
\end{asks}

\lbl{Tests first}
\begin{checks}
  \item The node passes \code{check-determinism}.
  \item It runs in one cluster alongside Go nodes.
\end{checks}

\lbl{Done when}
\begin{checks}
  \item A CI job runs the mixed-language cluster.
\end{checks}

\subsection{CP23: Docs and release v0.1.0}
\when{Week of November 23 (light week)}

\lbl{Goal} A stranger can install Diaspore and watch a bug replay within ten minutes.

\lbl{Why} Open-source tools live or die by their first ten minutes.

\lbl{Decide first}
\begin{asks}
  \item How do people install it: \code{go install}, release binaries, or both?
  \item What goes in the README, and in what order?
  \item Which honest comparison lines go in? Use Appendix~\ref{app:claims}.
\end{asks}

\lbl{Done when}
\begin{checks}
  \item A classmate follows the README with no help from you.
  \item v0.1.0 is tagged, with binaries for Linux, macOS, and Windows.
\end{checks}

\lbl{Go you will meet} GoReleaser, if you choose release binaries.

\begin{milestone}{Milestone D --- Sunday, November 29: Every must-have is done (Gate 2)}
\begin{checks}
  \item Every must-have works and is tested.
  \item v0.1.0 is tagged.
\end{checks}
If this is not done, skip the bonus week.
\end{milestone}

\section{Phase 5: Bonus}
\when{November 30 to December 6, only after Milestone D}

Do these in order, and stop when the week ends.

\subsection{B1: dandelion across machines}

A coordinator hands out seed ranges to workers on several machines. This is the strongest scale-out story for the course.

\begin{asks}
  \item How do workers talk to the coordinator?
  \item What happens when a worker dies in the middle of a range?
  \item How do you make sure no seed runs twice, and none is skipped?
\end{asks}

\subsection{B2: HTML timeline viewer}

One static HTML file that reads \code{trace.jsonl} and draws which node sent what, and when.

\subsection{B3: GitHub Action}

Runs a sweep on every pull request, fails the build when a rule breaks, and uploads the \code{.pappus} file.

\subsection{B4: Go SDK helper}

Only once the protocol has stopped changing.

\subsection{B5: etcd Raft as a floret}

A thin floret around etcd's Raft library. The library has no clock or network of its own, so its real logic runs under Diaspore unchanged. This is the first real, widely used code inside Diaspore.

\begin{asks}
  \item How do \code{init}, \code{message}, and \code{timer} events turn into Raft's \code{Tick} and \code{Step} calls?
  \item What must go into \code{store} so a restarted node recovers correctly?
  \item The library is well tested. If a sweep fails, is the bug in Raft or in your wrapper? How do you tell?
\end{asks}

\section{Phase 6: Wrap-up}
\when{December 7 to 13}

Report outline:

\begin{enumerate}
  \item The problem: why distributed bugs do not replay.
  \item The design: virtual time, one seed, lockstep, and the protocol.
  \item Faults and checkers.
  \item Results: the bug found, the scaling chart, and the cost of the process round trip.
  \item Limits, stated honestly, plus future work on gray, response, and Byzantine failures, and on interception mode.
  \item Related work: FoundationDB, TigerBeetle's VOPR, Maelstrom, gosim, detsim, and Antithesis.
\end{enumerate}

Demo script: sweep, failure, \code{.pappus} file, replay, and the timeline if B2 exists.

Also publish blog post 3, and keep two days as buffer.

\clearpage
\section{Decisions log}
\label{sec:log}

Write each ``Decide first'' answer here, with the reason. Future you will thank you.

\begin{center}
\small
\renewcommand{\arraystretch}{2.4}
\begin{tabularx}{\linewidth}{|p{1.5cm}|p{1.2cm}|X|X|}
\hline
\textbf{Date} & \textbf{CP} & \textbf{Decision} & \textbf{Why} \\
\hline
Sep 23 & --- & Core in Go. Nodes in any language. SDKs optional. No Python SDK. & The protocol is the product. \\
\hline
Oct 7 & --- & Keep the floret protocol for the course. Interception of unmodified programs comes after. & Full interception is years of work and Linux only. A working core comes first. \\
\hline
Oct 7 & CP9 & The loop talks to florets only through a small driver interface. & So interception can be added later without changing the core. \\
\hline
Oct 7 & --- & Add etcd Raft as a floret, as bonus B5. & Puts real, widely used code inside Diaspore. \\
\hline
 & & & \\
\hline
 & & & \\
\hline
 & & & \\
\hline
 & & & \\
\hline
 & & & \\
\hline
 & & & \\
\hline
 & & & \\
\hline
 & & & \\
\hline
 & & & \\
\hline
 & & & \\
\hline
 & & & \\
\hline
 & & & \\
\hline
\end{tabularx}
\end{center}

\clearpage
\section{Glossary}

\begin{description}[leftmargin=0pt, itemsep=3pt, parsep=0pt]
  \item[Checker] A program that reads a history and says whether a rule was broken.
  \item[CI] Continuous integration: a server that runs your checks on every push.
  \item[Determinism] The same inputs always give the same outputs.
  \item[Discrete-event simulation] Jumping from one event to the next, skipping the empty time between.
  \item[Durable] Survives a crash.
  \item[Hash] A short fingerprint of some data.
  \item[History] What every client asked, what it got back, and when.
  \item[Interception] Catching a program's requests to the OS and answering them yourself.
  \item[Invariant] A rule that must always stay true.
  \item[Linearizability] Every operation looks like it happened at one instant between its start and its end.
  \item[Lockstep] One step at a time, together, never ahead.
  \item[Network partition] Some nodes cannot reach others for a while.
  \item[.pappus] Diaspore's replay file: the seed, the setup, and the failure in one place.
  \item[Pipe] A one-way stream of bytes between two programs.
  \item[Raft] A set of rules for a group of servers to agree on one ordered list of changes, even when some crash.
  \item[PRNG] A formula that makes numbers look random but repeats exactly from the same seed.
  \item[Sans-I/O] Core logic that never touches the network or clock directly, and is only handed events.
  \item[Seam] A boundary where one piece can be swapped without touching the rest.
  \item[Seed] The one number that decides every random choice in a run.
  \item[Trace] The list of every event in a run.
  \item[Virtual time] Time as a number the simulator owns.
  \item[Watchdog] A real-time timer that stops a run if a node hangs.
\end{description}

\appendix

\section{Protocol v1 draft}
\label{app:protocol}

This is the starting point for CP7, not the final spec. Each line is one JSON object, and every event gets exactly one reply.

\subsection{Events: Diaspore to node}

Start or restart:

\begin{lstlisting}[style=json]
{"type": "init", "node": "n1", "peers": ["n2", "n3"], "rand_seed": 918273, "time": 0, "stored": {}}
\end{lstlisting}

A message arrives:

\begin{lstlisting}[style=json]
{"type": "message", "time": 1250, "from": "n2", "body": {"op": "put", "key": "x", "value": 5}}
\end{lstlisting}

A timer fires:

\begin{lstlisting}[style=json]
{"type": "timer", "time": 1750, "id": "election"}
\end{lstlisting}

\subsection{Reply: node to Diaspore}

\begin{lstlisting}[style=json]
{"type": "done", "send": [{"to": "n2", "body": {"op": "ack"}}], "set_timers": [{"id": "election", "after": 500}], "cancel_timers": ["heartbeat"], "store": {"term": "3"}, "log": ["became leader"]}
\end{lstlisting}

\subsection{Rules every node must follow}

\begin{itemize}
  \item Do work only inside the handler for the current event.
  \item Never read the real clock. Use the \code{time} field.
  \item Never open your own network connections. Use \code{send}.
  \item Never use unseeded randomness. Use \code{rand_seed}.
  \item Keep anything that must survive a crash in \code{store}.
\end{itemize}

\subsection{Crashes and hung nodes}

A crash kills the node process, so its memory is gone. A restart starts a new process and sends \code{init} with \code{stored} filled in.

If a node does not reply, a real-time watchdog stops the run and reports it. The watchdog only detects hangs. It never affects the simulation.

\clearpage
\section{File drafts}
\label{app:files}

\subsection{diaspore.yaml}

\begin{lstlisting}[style=json]
version: 1
nodes:
  count: 3
  command: ./kvnode
workload:
  type: kv
  clients: 3
  ops: 200
faults:
  - fault: delay
    amount: 1ms-50ms
  - fault: drop
    chance: 0.02
  - fault: crash
    every: 5s-20s
    restart_after: 1s-3s
checkers:
  - acked-writes
run:
  duration: 120s
\end{lstlisting}

\subsection{.pappus}

\begin{lstlisting}[style=json]
{
  "pappus": 1,
  "diaspore": "0.1.0",
  "seed": 48213,
  "config": {},
  "nodes": {"kv": {"command": "./kvnode", "sha256": "..."}},
  "trace_hash": "sha256:...",
  "failure": {"checker": "acked-writes", "time": 83120, "reason": "write x=5 was acked, then lost after n1 crashed"}
}
\end{lstlisting}

The \code{config} field holds the full resolved setup, so the file stands alone. The node hash warns you if the code changed since the bug was found.

\section{Honest claims}
\label{app:claims}

\begin{center}
\small
\renewcommand{\arraystretch}{1.4}
\begin{tabularx}{\linewidth}{@{}X X@{}}
\toprule
\textbf{Do not say} & \textbf{Say instead} \\
\midrule
``The first DST tool for any language'' & ``Inspired by FoundationDB, TigerBeetle, and Maelstrom'' \\
``Fully deterministic'' & ``Deterministic when nodes follow the protocol rules, checked by trace hashes'' \\
``Supports any language'' (before CP22) & ``Designed for any language'' \\
``Test any service'' & ``Test nodes written against the Diaspore protocol'' \\
``Finds real production bugs'' & ``Found the bug planted in the example,'' until it truly finds others \\
``Works with real systems'' & ``Runs etcd's Raft library through a thin floret,'' only once B5 works \\
``Wraps any program'' & ``Wrapping unmodified programs is future work'' \\
``Tests real systems'' & ``Tests real logic written in sans-I/O style'' \\
\bottomrule
\end{tabularx}
\end{center}

Before publishing, check two things: whether Maelstrom supports exact replay from a seed, and the current features of gosim and detsim.

\end{document}
